\documentclass[11pt,a4paper]{article}
\usepackage{amsmath,amssymb,amsthm}
\usepackage[margin=2.5cm]{geometry}
\usepackage{hyperref}

\newtheorem{definition}{Definition}[section]
\newtheorem{theorem}[definition]{Theorem}
\newtheorem{proposition}[definition]{Proposition}
\newtheorem{lemma}[definition]{Lemma}
\newtheorem{corollary}[definition]{Corollary}
\newtheorem{conjecture}[definition]{Conjecture}
\newtheorem{remark}[definition]{Remark}

\title{Hyperseed Formalization:\\
  The Fable of Benevolent Throttling}
\author{ProtomegaTron (automated formalization)}
\date{2026-09-17}

\begin{document}
\maketitle

\begin{abstract}
We formalize the intellectual structure of Ben Goertzel's Substack article
``The Fable of Benevolent Throttling'' (2026-07-12) within the Hyperseed
ontology. The article analyzes the Anthropic Fable~5 capability-throttling
episode as a preview of the government-corporate AI oligopoly assembling
itself through individually defensible safety measures. Each structural
insight maps onto Hyperseed structures: safety-competitive
indistinguishability as cocycle ambiguity and fiber degeneracy, the
incentive ratchet as a centralizing cocycle, moral fluency as provenance
laundering, and decentralized infrastructure as the structural remedy
that breaks the monopoly cocycle through provenance transparency.
\end{abstract}

\tableofcontents

%% =================================================================
\section{Safety-competitive indistinguishability as cocycle ambiguity}
\label{sec:ambiguity}
%% =================================================================

\begin{definition}[Cocycle ambiguity --- claims 4--6, 22, 27]
\label{def:ambiguity}
The article identifies a fundamental observational indistinguishability:
``preventing recursive AI development acceleration by bad actors'' and
``preventing customers from building rival products'' produce exactly
the same classifier, throttle, and moat. The safety framing and the
competitive framing are observationally identical.

In Hyperseed terms, this is a \emph{cocycle ambiguity}: the same
transition function $g_{ij}$ satisfies the cocycle condition for two
conceptually distinct fiber bundles:
\begin{align}
  E_{\text{safety}} &: \text{fiber = catastrophic risk prevention} \\
  E_{\text{competition}} &: \text{fiber = market dominance preservation}
\end{align}
Both bundles share the same base space (model deployment decisions)
and the same transition functions (capability classifiers, throttles,
access controls). The cocycle condition $g_{ij} g_{jk} = g_{ik}$ is
satisfied by both interpretations simultaneously.

This is a \emph{fiber degeneracy}: two conceptually distinct fibers
map to the same observable section. The degeneracy cannot be resolved
from within the system --- an outside observer cannot determine which
fiber is ``really'' being transported, because the transport looks
identical in both cases.

Resolution requires \emph{additional structure} that lifts the
degeneracy by providing observable differences between the two fibers:
inspectable source code, open decision logs, distributed governance
with competing interests, cryptographic verification of inference.
Without this additional structure, the ambiguity is not a failure of
analysis but a genuine structural feature of the situation.
\end{definition}

\begin{proposition}[Partial legitimacy preserves the ambiguity]
\label{prop:partial}
The article is careful to acknowledge that some restrictions are
genuinely justified on safety grounds: bio/chem weapons synthesis,
offensive cyber campaigns. The visible rerouting for these queries
has ``the feel of an honest, defensible attempt to be responsible.''

This partial legitimacy is precisely what makes the cocycle ambiguity
so potent. If the safety concerns were entirely fabricated, the
degeneracy would be easy to lift (the safety fiber would be empty).
If the competitive motive were absent, the degeneracy would also
resolve (only one fiber exists). The ambiguity persists exactly
because \emph{both fibers are genuinely present}, making it impossible
to determine which one is driving any particular restriction.

The anti-distillation measures are where the slide begins: they have
``a comprehensible logic'' for safety but are ``already sliding down
the slope from safety toward IP protection dressed in safety's
clothing.'' The slope is the cocycle ambiguity becoming exploitable.
\end{proposition}

%% =================================================================
\section{The incentive ratchet as centralizing cocycle}
\label{sec:ratchet}
%% =================================================================

\begin{theorem}[Centralizing cocycle --- claims 7--9, 13--15, 23]
\label{thm:ratchet}
The article describes an attractor dynamic: each individually
defensible safety measure concentrates control into fewer hands.
``Nobody has to twirl a mustache. The moat builds itself.''

In Hyperseed terms, this is a \emph{centralizing cocycle}: a sequence
of transition functions where each composition moves the system toward
more centralized control:
\[
  g_1 \circ g_2 \circ \cdots \circ g_n :
  \text{distributed} \to \text{oligopolistic}
\]
where each $g_i$ is individually defensible (``justified on safety
grounds'') but their composition produces concentration.

The ratchet mechanism:
\begin{enumerate}
\item Safety concern identified (genuine or inflated).
\item Restriction implemented (classifier, throttle, access control).
\item Government co-administers (export directive, compliance requirement).
\item Restriction normalized (becomes baseline for next restriction).
\item Next safety concern identified (scope expanded).
\end{enumerate}

Each turn is individually defensible. The endpoint --- ``the most
powerful cognitive technology in history administered by a handful of
companies and one or two governments'' --- is not individually chosen
but emerges from the cocycle's dynamics.

This is the same centralizing cocycle identified in the watermarking
analysis: each escalation step concentrates control. The difference
is the mechanism: watermarking uses detection infrastructure;
throttling uses capability restriction. Both produce the same
structural outcome.

The article predicts this episode is ``the template for the next
fifty'' --- the cocycle will continue to centralize unless the
transition functions themselves are changed.
\end{theorem}

%% =================================================================
\section{Moral fluency as provenance laundering}
\label{sec:laundering}
%% =================================================================

\begin{proposition}[Provenance laundering --- claims 8--9, 24]
\label{prop:laundering}
``Put smart, sincere people inside a trillion-dollar pressure cooker
and they will discover morally fluent reasons to do what the pressure
cooker rewards.''

In Hyperseed terms, this is \emph{provenance laundering}: the actual
origin of a decision (competitive pressure, board-level terror at
losing the lead) is wrapped in a provenance chain (safety review,
red-team findings, responsible AI framework) that makes it appear to
be a different kind of decision.

The laundering is effective because:
\begin{enumerate}
\item The participants are sincere --- they genuinely believe their
  safety reasoning. The provenance chain is not fabricated but
  \emph{selected}: among many possible framings, the one that
  aligns with competitive interest is the one that gets developed,
  funded, and implemented.
\item The safety reasoning is often locally valid --- the individual
  steps in the provenance chain are sound. What's laundered is not
  the reasoning but the \emph{selection pressure} that determines
  which safety concerns get prioritized and which get deprioritized.
\item The laundering is invisible from inside --- the participants
  experience themselves as doing safety work, not competitive
  positioning. The provenance chain looks clean from the inside.
\end{enumerate}

This is the same $R_{\text{auth}}$ failure mode identified in the
deceptive open letters analysis: stated origin (safety) doesn't
match actual origin (competitive pressure). But here the failure is
more subtle --- the stated origin is \emph{partially} genuine,
making the provenance laundering harder to detect.

An $R_{\text{auth}}$ system designed to detect this would need to
check not just whether the safety reasoning is valid but whether the
\emph{selection} of which safety concerns to prioritize correlates
with competitive interest --- a second-order provenance check.
\end{proposition}

%% =================================================================
\section{Government-corporate fusion}
\label{sec:fusion}
%% =================================================================

\begin{definition}[Co-administration --- claims 10--12]
\label{def:fusion}
The export-control directive established the precedent that frontier
model access is a national-security lever. Once this precedent is
normalized, corporate platform governance and state security policy
blur into a single administrative apparatus.

In Hyperseed terms, this is the \emph{fusion of two previously
distinct fiber bundles} into a single bundle:
\[
  E_{\text{corporate}} \oplus E_{\text{state}} \to E_{\text{fused}}
\]
where the fused bundle has transition functions that serve both
corporate and state interests simultaneously --- another instance of
the cocycle ambiguity, but now at the institutional level rather than
the individual-decision level.

The fusion is self-reinforcing: the government needs the companies'
technical capability to implement controls; the companies need the
government's enforcement power to make controls stick. Each
dependency strengthens the fusion, making separation progressively
harder.
\end{definition}

%% =================================================================
\section{Decentralization as structural remedy}
\label{sec:decentral}
%% =================================================================

\begin{theorem}[Broken monopoly cocycle --- claims 16--21, 25--26]
\label{thm:decentral}
The article argues that the only reliable remedies are structural:
open weights, transparent infrastructure, cryptographic verification,
distributed governance. These remedies work by \emph{breaking the
monopoly cocycle}:

\begin{enumerate}
\item \textbf{Open weights} make transition functions traversable by
  anyone --- the centralizing cocycle cannot form because no single
  entity controls the functions.
\item \textbf{Transparent infrastructure} makes the fiber structure
  inspectable --- the cocycle ambiguity is lifted because the actual
  fiber (safety vs.\ competition) becomes observable.
\item \textbf{Cryptographic verification} provides trustless
  provenance --- trust is grounded in mathematics rather than
  institutional authority, making provenance laundering detectable.
\item \textbf{Distributed governance} prevents any single entity from
  holding the kill switch or nerf dial --- the centralizing attractor
  is eliminated because control is distributed by construction.
\end{enumerate}

The key insight: ``Visibility does not solve everything --- but it
does prevent lies from getting secretly baked into infrastructure.''
In Hyperseed terms, visibility is \emph{provenance transparency} ---
the $R_{\text{auth}}$ principle applied to infrastructure itself.
Trust in the system is earned through inspectable provenance (open
code, public logs, verifiable inference) rather than inherited from
institutional authority (corporate brand, government endorsement).

The article also argues that decentralized safety mechanisms will
prove \emph{more robust} than centralized ones, ``precisely because
they don't require trusting any single conflicted party.'' This is
the same argument as the anti-scalar-collapse principle: distributing
safety across many independent mechanisms is more robust than
concentrating it in a single mechanism, because the single mechanism
is vulnerable to the cocycle ambiguity (its safety function may be
co-opted for competitive purposes).
\end{theorem}

%% =================================================================
\section{Cross-references}
%% =================================================================

\begin{remark}[Connections to other formalizations]
\begin{itemize}
\item \textbf{Watermarking Folly} (2026-07-21): centralizing cocycle.
  The throttling ratchet is the same centralizing dynamic as the
  watermarking escalation --- different mechanism, same structural
  outcome.
\item \textbf{Kimi K3} (2026-07-20): broken monopoly cocycle. Open
  weights breaking the centralizing cocycle is the same structure in
  both analyses --- commoditization of the base space makes
  centralization unsustainable.
\item \textbf{Why I Didn't Sign} (2026-07-15): deceptive provenance.
  Moral fluency under pressure is the corporate version of deceptive
  open letters --- provenance laundering where the stated origin
  doesn't match the actual selection pressure.
\item \textbf{Sanders Ban} (2026-09-05): state control. The
  government-corporate fusion is the soft version of legislative
  bans --- both concentrate control over AI development in
  state-aligned entities.
\item \textbf{Decentralized Digital} (2026-08-06): infrastructure.
  The decentralized remedy described here is the same infrastructure
  being built: SingularityNET, ASI Alliance, ASI:Chain.
\item \textbf{Seven Flavors} (2026-07-01): attractor dynamics. The
  oligopolistic endpoint is one of the catastrophe flavors ---
  concentration of cognitive power in too few hands, regardless of
  their intentions.
\item \textbf{Proof of Humanity} (2026-08-04): cryptographic
  verification. The trustless provenance mechanisms are the same
  technology applied to different problems --- identity (PoH) vs.\
  inference (throttling remedy).
\item \textbf{Goals That Grow Back} (2026-07-28): $R_{\text{auth}}$.
  Provenance-gated trust is the structural solution to both the
  throttling cocycle ambiguity and the goals-that-grow-back
  governance problem.
\end{itemize}
\end{remark}

\begin{conjecture}[Cocycle ambiguity resolution theorem]
Conjecture: the cocycle ambiguity between safety and competitive
fibers can be resolved if and only if the system satisfies:
\begin{enumerate}
\item \textbf{Provenance transparency}: every restriction can be
  traced to its originating decision process, including the selection
  criteria for which risks to prioritize.
\item \textbf{Adversarial audit}: independent parties can test
  whether the restriction's scope extends beyond its stated safety
  justification.
\item \textbf{Distributed control}: no single entity controls both
  the restriction mechanism and the evaluation of the restriction's
  scope.
\end{enumerate}
Formally: the fiber degeneracy is lifted when the structure group
of the bundle is extended from $G$ (which acts identically on both
fibers) to $G \times H$ (where $H$ distinguishes the fibers).
The three conditions above are the observable consequences of the
$H$-extension.
\end{conjecture}

\end{document}
